% Contenido exclusivo del propietario de entropías tipadas.
\subsection{Correlación cuasilibre y entropías espectrales de ocupación}

Para un estado fermiónico cuasilibre gauge-invariante, conservador del
número y definido sobre un espacio de un cuerpo de dimensión finita, sea
\(C_{ij}=\langle a_j^\dagger a_i\rangle\) el operador de correlación de un
cuerpo. La hipótesis gauge-invariante excluye correlaciones anómalas de
emparejamiento, de modo que \(C\) determina el estado cuasilibre. Si
\(0<C<I\), su generador modular de un cuerpo y la reconstrucción recíproca
del correlador son
\begin{equation}
 \boxed{
 K_1=\log\bigl((I-C)C^{-1}\bigr),
 \qquad
 C=(I+e^{K_1})^{-1}.}
 \label{eq:puente-cuasilibre-modular-rev3}
\end{equation}

\begin{corollary}[Forma logística del espectro modular]
\label{cor:forma-logistica-modular-rev3}
Si \(K_1u_i=\kappa_i u_i\), la ocupación del modo \(u_i\) es
\[
 \langle n_i\rangle=\frac1{1+e^{\kappa_i}}.
\]
Para \(\kappa_i=\beta(\varepsilon_i-\mu)\), se recupera la ley de
Fermi--Dirac.
\end{corollary}

\begin{proof}
La diagonalización espectral de \(K_1\) diagonaliza toda función suya. La
aplicación de \(x\mapsto(1+e^x)^{-1}\) a cada autovalor produce la fórmula.
\end{proof}

Si \(C\) tiene autovalores \(p_i\in[0,1]\), la entropía de ocupación del
estado fermiónico cuasilibre es, con la convención \(0\log0=0\),
\begin{equation}
 S_F=-\sum_i\bigl[p_i\log p_i+(1-p_i)\log(1-p_i)\bigr].
 \label{eq:entropia-cuasilibre-fermi-rev3}
\end{equation}
Para un estado bosónico gaussiano diagonal con un número finito de modos y
ocupaciones medias \(\overline n_i\),
\begin{equation}
 S_B=\sum_i\bigl[(1+\overline n_i)\log(1+\overline n_i)
                  -\overline n_i\log\overline n_i\bigr].
 \label{eq:entropia-gaussiana-bose-rev3}
\end{equation}
La misma fórmula se extiende a una familia numerable cuando la serie del
miembro derecho converge.
Estas entropías cuantifican configuraciones de ocupación. La entropía de
entrelazamiento posee otro dominio: el estado reducido inducido por una
descomposición tensorial. La identificación entre ambas requiere declarar la
reducción y el operador de correlación que la implementan.

Finalmente, para un estado reducido \(\rho_A\), la \emph{entropía de von
Neumann}
\[
 S_{\rm vN}(\rho_A)=-\operatorname{Tr}(\rho_A\log\rho_A)
\]
es una propiedad espectral del operador de densidad. Cuando \(\rho_A\)
procede de una purificación global, cuantifica el entrelazamiento a través
del corte. En la construcción de borde HMT, el operador de área, el
Hamiltoniano modular y la reducción correspondiente permanecen en la
\cref{hap:sec:area-entrelazamiento-barbero-rev11}.

Los mapas entre estos dominios se publican mediante identidades explícitas.
La descomposición de una medida por las fibras de \(q\) satisface la regla de
la cadena
\begin{equation}
 \mathsf H_\mu(\mathbf X)
 =\mathsf H_{q_*\mu}(q(\mathbf X))
  +\mathsf H_{\rm fib}(q;\mu).
 \label{eq:cadena-entropia-fibras-rev2}
\end{equation}
\Needspace{5\baselineskip}
Una distribución finita \(p=(p_i)\) se realiza diagonalmente como
\[
 \rho_p=\sum_i p_i|i\rangle\langle i|,
 \qquad S_{\rm vN}(\rho_p)=\mathsf H(p).
\]
Si \(P_{N,E}\) proyecta sobre el subespacio microcanónico de dimensión
\(\Omega(N,E)\), entonces
\[
 \rho_{\rm mc}=\frac{P_{N,E}}{\Omega(N,E)},
 \qquad S_{\rm vN}(\rho_{\rm mc})=\log\Omega(N,E)=S_{\rm mc}.
\]
Para un estado puro bipartito,
\[
 \rho_A=\operatorname{Tr}_B
 \bigl(|\Psi_{AB}\rangle\langle\Psi_{AB}|\bigr)
\]
define la reducción que mide el entrelazamiento. Finalmente, la igualdad entre
\(h_{\rm top}=\log\rho(A)\) y una entropía métrica se alcanza, para este
sistema primitivo, mediante la medida de Parry, que realiza la entropía métrica
máxima; no se atribuye a una medida arbitraria de historias.

